\section{Where B-splines, Gaussian mixtures and neural networks fit}\label{sec-alternatives}

\subsection{What B-spline coefficients cannot see}
\begin{proposition}[Leverage of a spike]\label{prop-leverage}
Let $B\in\R^{K\times q}$ be the B-spline design matrix at the pillars and $H=B(B^\top B)^{-1}B^\top$ the hat matrix. The
least-squares fit of a unit spike $e_k$ has energy $\|He_k\|^2=h_{kk}$, and $\sum_kh_{kk}=q$. A fraction $1-h_{kk}$ of the
spike's energy stays in the residual and is invisible to any feature computed from the coefficients.
\end{proposition}
\begin{proof}
$H$ is symmetric and idempotent, so $\|He_k\|^2=e_k^\top H^\top He_k=e_k^\top He_k=h_{kk}$, and
$\sum_kh_{kk}=\operatorname{tr}H=\rank B=q$. The coefficients $(B^\top B)^{-1}B^\top e_k$ determine $He_k$ and nothing else.
\end{proof}

\begin{example}
For the six cubic B-splines (knots in $\ln T$) of the benchmark and the eleven pillars,
\begin{center}\input{generated/\lang/tab_leverage}\end{center}
At the interior pillars between nine months and four years the coefficients capture only \vLevMin{} to \vLevMax{} of a
spike's energy (on average $6/11\approx\num{0.545}$ over all pillars). The spline also spreads what it captures over
several coefficients. A one-pillar anomaly, the most common bad-curve pattern, is therefore diluted before any
clustering starts.
\end{example}

\subsection{Why a mixture on the difference flags FO mismarks}
\begin{proposition}\label{prop-gmm-fo}
Let the features be the spline coefficients $c=B^+\Delta b$ of the repo difference. A parallel FO repo shift $\delta$
moves them by $-\delta\,B^+\one$, a fixed direction. If the mixture is fitted on clean history, where FO and Totem agree
up to noise, its density decreases along that direction, so its alarm rate on FO mismarks tends to one as $|\delta|$
grows, whatever the quality of the Totem curve.
\end{proposition}
\begin{proof}
$c$ is linear in $\Delta b$, and $\Delta b$ moves by $-\delta\one$. Each Gaussian component has a finite variance along
$B^+\one$, so the mixture density, and hence the likelihood score, goes to zero along that direction.
\end{proof}
The structural detector has no such problem: $T\cdot\one$ lies in the span of the null model, so a level difference is
projected out of every test and reported through $\hat\beta_0$. The benchmark confirms both statements: \vFoGmmDiff{} of
FO mismarks are flagged by the mixture on the difference, \vFoStruct{} by the structural detector.

\subsection{Where they help}
\begin{itemize}
\item \textbf{P-splines in the null model.} If $\beta_0+\beta_1\ln T$ proves too rigid on real data, a penalised spline
\cite{eilers1996} can replace it as the smooth part of the null model; the tests of \cref{sec-detection} are unchanged.
\item \textbf{Mixtures on structured outputs.} A Gaussian mixture on the detector's outputs ($\hat\beta_0$, $\hat\beta_1$, the
$z_t$, $Q$ and the dispersion ratios) learns the usual regimes per region and category in a low-dimensional space where each coordinate means something.
\item \textbf{Region $\times$ category} is best used to group the calibration of $\kappa$ and of the level scale, not as
clusters.
\item \textbf{Neural networks.} By \cref{prop-nn}, on the mechanisms we can name a network can at best learn the
matched filter. Networks earn their place once there is a labelled incident history and patterns the templates do not
cover; they should then be trained on residuals of the structural model and on synthetic injections, not on raw
curves.
\end{itemize}
